\section{\Kcas Overview}

\subsection{Shared-Memory Locations}

\begin{figure}
\begin{ocamlcode}
module Loc : sig
  type !'a t
  
  val make : ?padded:bool -> ?mode:Mode.t -> 'a -> 'a t
  val make_array : ?padded:bool -> ?mode:Mode.t -> int -> 'a -> 'a t array
  
  val get : 'a t -> 'a
  val set : 'a t -> 'a -> unit
  val exchange : 'a t -> 'a -> 'a
  val compare_and_set : 'a t -> 'a -> 'a -> bool
  val fetch_and_add : int t -> int -> int
  val incr : int t -> unit
  val decr : int t -> unit
  
  (* [get_as f l] is equivalent to [f (get l)] except [f] may raise the
     exception [Retry.Later] to retry after [l] has been modified. *)
  val get_as : ?timeoutf:float -> ('a -> 'b) -> 'a t -> 'b
  (* [update l f] repeats [let x = get l in compare_and_set l x (f x)] until it
     succeeds. [f] may raise the exception [Retry.Later] to retry after [l] has
     been modified. *)
  val update : ?timeoutf:float -> 'a t -> ('a -> 'a) -> 'a
  (* [modify l f] is equivalent to [update l f |> ignore]. *)
  val modify : ?timeoutf:float -> 'a t -> ('a -> 'a) -> unit
end
\end{ocamlcode}
\caption{Interface for shared-memory locations}
\label{fig:loc_interface}
\end{figure}

\subsection{Transaction Log}

\begin{figure}
\begin{ocamlcode}
module Xt : sig
  type 'x t
  
  val get : xt:'x t -> 'a Loc.t -> 'a
  val set : xt:'x t -> 'a Loc.t -> 'a -> unit
  val exchange : xt:'x t -> 'a Loc.t -> 'a -> 'a
  val compare_and_set : xt:'x t -> 'a Loc.t -> 'a -> 'a -> bool
  val fetch_and_add : xt:'x t -> int Loc.t -> int -> int
  val incr : xt:'x t -> int Loc.t -> unit
  val decr : xt:'x t -> int Loc.t -> unit
  val swap : xt:'x t -> 'a Loc.t -> 'a Loc.t -> unit
  
  val update : xt:'x t -> 'a Loc.t -> ('a -> 'a) -> 'a
  val modify : xt:'x t -> 'a Loc.t -> ('a -> 'a) -> unit
  
  type 'a tx = { tx : 'x. xt:'x t -> 'a } [@@unboxed]
  
  val commit : ?timeoutf:float -> ?mode:Mode.t -> 'a tx -> 'a
  val post_commit : xt:'x t -> (unit -> unit) -> unit
  
  val to_blocking : xt:'x t -> (xt:'x t -> 'a option) -> 'a
  val to_nonblocking : xt:'x t -> (xt:'x t -> 'a) -> 'a option
  
  val first : xt:'x t -> (xt:'x t -> 'a) list -> 'a
  
  type 'x snap
  val snapshot : xt:'x t -> 'x snap
  val rollback : xt:'x t -> 'x snap -> unit
  
  val validate : xt:'x t -> 'a Loc.t -> unit
end
\end{ocamlcode}
\caption{Interface for explicit transaction logs}
\label{fig:xt_interface}
\end{figure}

transactional API

A transaction is a specification for generating a list of CAS

\begin{ocamlcode}
let transfer ~xt () =
  let a' = Xt.get ~xt a
  and b' = Xt.get ~xt b in
  Xt.set ~xt s (a' + b')
\end{ocamlcode}

\begin{ocamlcode}
Xt.commit { tx = transfer }
\end{ocamlcode}

\subsection{Retry Mechanism}

% retry mechanism (wait until a read variable is modified) (as in Composable Memory Transactions)
% see [Xt.first]
% [retry] uses the transaction log generated by the failed transaction to determine which memory cells it read, and automatically retries the transaction when one of these cells is modified
% [retry] can, at any time, abort the transaction and wait until some value previously read by the transaction is modified before retrying

\subsection{Read Skews}

% read skews / conflicts :
% One of the variable read by a transaction X is modified by another transaction Y during the execution of X
% Such a transaction is doomed to abort if it ever tries to commit, but it can cause issues since invariants can appear to be broken
% solution :
% check the conflict only at the end of the transaction
% Avoid instrumenting all the reads or all the writes
% Possible inconsistency if a transaction continues to run with invalid values (typically in case of double read conflicts)
% in Transactional Locking II (using a global clock), the implementation guarantees that the code is executed on a consistent snapshot of memory
% abort any transaction that is not valid

\subsection{Standard Concurrent Data Structures}